% Quarantine-only candidate proof.
% Date: 29 August 2026.
% This file is not a live-book input and makes no novelty or priority claim.

\begin{proof}[Proof of
\cref{lem:p213-normal-curve-nonrigid-localization}]
Because $Y$ is good and separated, after shrinking inside $W$ we may
choose an affinoid analytic neighborhood of $y$.  Its connected component
containing $y$ is again an affinoid analytic neighborhood: the connected
components of an affinoid space are the finitely many clopen affinoid
domains cut out by its idempotents.  Denote this component by $Q$.  It is
compact, connected, and holomorphically separated; the last assertion is
part of the affinoid case of \cite[\S15.2.2]{CLD25}.  Analytic domains of a
normal analytic space are normal.  We record why connectedness then makes
$Q$ irreducible.  If two irreducible components of $Q$ met at a point
$q$, the local ring $\mathcal O_{Q,q}$ would have at least two minimal
prime ideals and hence would not be a domain.  This is impossible because
$Q$ is normal.  Thus the irreducible components are pairwise disjoint.
There are only finitely many of them because $Q$ is affinoid, so each is
both closed and open.  The connectedness of $Q$ now shows that there is
only one.  The point $y$ lies in the relative interior of $Q$ in $W$.

If $\partial(Q/W)$ is empty, take $V=Q$ and take for $F$ the constant unit
moment $1:Q\to\mathbf G_{m,K}$.  Then the asserted boundary avoidance is
immediate.  For an analytic-domain inclusion the relative boundary is the
ordinary topological boundary by \cite[\S11.4.1]{CLD25}; hence $Q$ is both
open and compact in $W$, so the final properness assertion also follows.
We may therefore suppose from now on that $\partial(Q/W)$ is nonempty.

We first isolate the point at which the proof of
\cite[Lemma~12.3.6]{CLD25} used a nontrivial value group.  Fix
$z\in\partial(Q/W)$.  Holomorphic separation gives
$f_z\in\mathcal O(Q)$ such that
\begin{equation}
\label{eq:p213-normal-curve-separator}
                         |f_z(z)|\ne|f_z(y)|.
\end{equation}
We claim that $f_z(y)\ne0$.  Suppose otherwise.  Since $Y$ is normal and
$y$ is nonrigid, the local description of normal analytic curves in
\cite[\S4.2.18]{Ducros24CurveStructure} gives
\[
                         \mathcal O_{Y,y}=\kappa(y).
\]
Thus the germ of $f_z$ at $y$ is zero, so $f_z$ vanishes on a nonempty
ordinary open subset of $Q$.  If $f_z$ were not identically zero, its zero
locus would be a proper Zariski closed subset of the irreducible
one-dimensional space $Q$.  Such a subset has dimension zero by
\cite[\S11.2.5]{CLD25}, whereas a nonempty analytic open subset of the pure
curve $Q$ has dimension one by \cite[\S11.2.3]{CLD25}.
This is impossible.  Hence $f_z$ is identically zero on $Q$, contradicting
\eqref{eq:p213-normal-curve-separator}.  The claim follows.

For each $z\in\partial(Q/W)$, put
\[
 Q_z:=\{x\in Q:|f_z(x)|\ne|f_z(y)|\}.
\]
These open sets cover the compact relative boundary.  Choose
$z_1,\ldots,z_N$ such that $Q_{z_1},\ldots,Q_{z_N}$ still cover it, and
write $f_i=f_{z_i}$.  Set
\[
                 r:=\frac12\min_{1\le i\le N}|f_i(y)|>0
\]
and let
\[
                 V:=\{x\in Q:|f_i(x)|\ge r
                              \text{ for }1\le i\le N\}.
\]
This is a compact analytic domain in $Q$ and an analytic neighborhood of
$y$.  Every $f_i$ is invertible on $V$, so the tuple
$F=(f_1,\ldots,f_N)$ is a moment on $V$.

For $0\leq j\leq N$, put
\[
 V_0:=Q,
 \qquad
 V_j:=\{x\in V_{j-1}:|f_j(x)|\geq r\}.
\]
Thus $V_N=V$.  The composition rule for relative interiors and the
boundary computation for domains defined by one inequality in
\cite[\S\S11.4.1--11.4.2]{CLD25}, applied to each successive
analytic-domain inclusion, give
\[
 \partial(V_j/W)
 \subseteq
 \bigl(\partial(V_{j-1}/W)\cap V_j\bigr)
 \cup\{x\in V_j:|f_j(x)|=r\}.
\]
Induction on $j$ therefore yields
\[
 \partial(V/W)
 \subseteq
 \bigl(\partial(Q/W)\cap V\bigr)
 \cup\bigcup_{i=1}^N\{x\in V:|f_i(x)|=r\}.
\]
Let $x\in\partial(V/W)$.  If $|f_i(x)|=r$ for some $i$, then
$|f_i(x)|\ne|f_i(y)|$ by the definition of $r$.  Otherwise
$x\in\partial(Q/W)$, and the finite covering of this boundary supplies an
$i$ such that $x\in Q_{z_i}$; again
$|f_i(x)|\ne|f_i(y)|$.  Thus
$F_{\mathrm{trop}}(x)\ne F_{\mathrm{trop}}(y)$ in both cases.  This proves
\eqref{eq:p213-normal-curve-nonrigid-boundary}.

The last assertion is the same localization argument as in the proof of
\cite[Lemma~12.3.6 and Theorem~16.8.5]{CLD25}.  The set $\Omega$ is open
because the relative boundary is compact.  Its inverse image $U$ misses
$\partial(V/W)$, hence lies in the relative interior of $V$ and is open in
$W$; it contains $y$ by
\eqref{eq:p213-normal-curve-nonrigid-boundary}.  Finally,
$F_{\mathrm{trop}}|_U$ is obtained from the proper map
$F_{\mathrm{trop}}:V\to\mathbf R^N$ by base change along the open immersion
$\Omega\hookrightarrow\mathbf R^N$, and is therefore proper.
\end{proof}
